\documentclass[11pt]{article}

% --------- BASIC PACKAGES ----------
\usepackage[left=0.55in,top=0.5in,right=0.55in,bottom=0.5in]{geometry}
\usepackage{enumitem}
\usepackage[hidelinks]{hyperref}
\usepackage{titlesec}
\usepackage{array}
\usepackage{setspace}

\usepackage[T1]{fontenc}
\usepackage{mathptmx} % SANDBOX SUBSTITUTE -- restore \usepackage{newtxtext,newtxmath} locally
\input{glyphtounicode}
\pdfgentounicode=1

\setstretch{0.97}
\pagenumbering{gobble}

% % ---------- SECTION FORMATTING ----------
% \titleformat{\section}{\large\scshape}{}{4em}{}[\vspace{-0.7em}\rule{\linewidth}{0.4pt}\vspace{-0.4em}]
% \titlespacing*{\section}{0pt}{0.45em}{0.25em}

% \newenvironment{cvrole}[5]{%
%   \par\noindent\textbf{#1} | \textit{\small #2, #3}\hfill \textit{\small #4 -- #5}\par
%   \begin{itemize}[leftmargin=1.5em, itemsep=-0.35em, topsep=-0.4em]%
% }{%
%   \end{itemize}\par\noindent\mbox{}\par\vspace{-0.45em}%
% }


% ---------- SECTION FORMATTING ----------
\titleformat{\section}{\Large\scshape}{}{4em}{}[\vspace{-0.7em}\rule{\linewidth}{0.4pt}\vspace{-0.4em}]
\titlespacing*{\section}{0pt}{1em}{0.4em}

\newenvironment{cvrole}[5]{%
  \par\noindent\textbf{\large #1} | \textit{\small #2, #3}\hfill \textit{\small #4 -- #5}\par
  \begin{itemize}[leftmargin=2em, itemsep=-0.3em, topsep=-0.4em]%
}{%
  \end{itemize}\par\noindent\mbox{}\par\vspace{-0.1em}%
}


% ---------- HEADER ----------
\begin{document}
\pagestyle{empty}

\begin{center}
    {\LARGE \scshape \textbf{Nishal Stanislaus Silva}}{\large \textit {, PhD}}
    \\[1pt]
    {\small Machine Learning Researcher | Applied AI, Evaluation \& Benchmark Design, Production ML Systems}
    \\[1pt]
    {\footnotesize\textit{Canadian Permanent Resident, Eligible to work in Canada without sponsorship}}
    \\[1pt]
    {\small
    \href{mailto:nishal.silva@hotmail.com}{nishal.silva@hotmail.com}
    \,\,|\,\, +1 613 200 2030
    \,\,|\,\, Toronto, ON (open to relocation)
    \,\,|\,\, \href{https://nishal.xyz}{nishal.xyz}
    \,\,|\,\, \href{https://www.linkedin.com/in/nishal-silva}{linkedin.com/in/nishal-silva}}
\end{center}

% =====================================================
\section*{Professional Summary}

Machine learning researcher with a PhD in Information Engineering and Computer Science and 10+ years of applied research and production experience. Experienced in evaluation: designing benchmarks, building datasets, and constructing ablation and regression infrastructures. I have released four public benchmark datasets and ten peer-reviewed papers, and deployed into regulated production settings, including manufacturing,  financial forensics, and compliance monitoring. 

% =====================================================
\section*{Core Competencies}

\textbf{Evaluation \& Benchmarking:} Benchmark Construction, Golden-Set Dataset Design, Ablation Studies, Metric Selection, Offline Evaluation Design, Regression Testing, Reproducible Pipelines\\
\textbf{Applied Machine Learning:} End-to-End ML Lifecycle, Deep Learning, Sequence Modeling, Retrieval \& Similarity Search, Generative Modeling, Multimodal Modeling, Anomaly Detection, Human-in-the-Loop Systems\\
\textbf{Systems \& Domain:} Production Deployment, Distributed Event-Driven Architectures, Financial Forensics \& Compliance Monitoring, Multi-Agent Coordination, Auditability, Low-Latency Inference

% =====================================================
\section*{Technical Stack}

\textbf{Languages:} Python, SQL, C++, C, MATLAB, JavaScript, Linux Shell Scripting\\
\textbf{Machine Learning:} PyTorch, TensorFlow, Keras, scikit-learn, HuggingFace, JAX, ONNX, NumPy, Pandas, OpenCV\\
\textbf{LLM \& Agentic:} Commercial LLM APIs (Anthropic, OpenAI), Retrieval-Augmented Generation, Prompt Engineering, Structured Output Generation\\
\textbf{Data \& Cloud:} PostgreSQL, Snowflake, AWS (EC2, S3, ECS, MWAA, RDS), ETL Pipelines, Information Retrieval\\
\textbf{Development \& MLOps:} Git, Docker, Linux/UNIX, Jenkins, CI/CD, Apache Airflow, Pytest, MLflow, REST APIs

% =====================================================
\section*{Education}
\textbf{Ph.D - Information Engineering and Computer Science} \textit{\small{ - University of Trento, Italy}} \hfill \textit{Jan 2025}\\[1pt]
\noindent\textbf{M.Sc - Telecommunication and Electronic Engineering} \textit{\small{ - Sheffield Hallam University, UK}} \hfill \textit{Aug 2018}\\[1pt]
\noindent\textbf{B.Eng (Hons) - Electronic Engineering} \textit{\small{ - Sheffield Hallam University, UK}} \hfill \textit{Mar 2014}

% =====================================================
\section*{Professional Experience}

\begin{cvrole}{Independent Machine Learning Researcher}{Self-directed}{Toronto, Canada}{Jan 2026}{Present}
    \item Researching foundation model based music generation and controllable synthetic variation generation, working with transformer and diffusion architectures and open-weight checkpoints.
    \item Maintain Nebula, an open-source C++ audio feature extraction library, and four public benchmark datasets on Zenodo used for evaluating real-time pattern detection systems.
\end{cvrole}

\begin{cvrole}{Postdoctoral Researcher}{University of Trento}{Trento, Italy}{Jan 2025}{Dec 2025}
    \item Built evaluation infrastructures within the EU-funded MUSMET project: golden-set benchmarks, baseline comparisons, ablation studies, and reproducible pipelines.
    \item Designed and released public benchmark datasets with annotation protocols and held-out evaluation splits.
    \item Designed Hot Licks Mapper: a distributed, event-driven system where on-device models autonomously coordinated a network of peripheral devices over custom protocols, with strict latency budgets and failure handling.
    \item Achieved 14\,ms inference at F1 0.76 on Raspberry Pi 4, a 74$\times$ speedup over deterministic baselines (JAES 2026); translated prototypes into documented components with academic and industry partners.
\end{cvrole}

\begin{cvrole}{Visiting Researcher}{McGill University}{Montreal, Canada}{May 2024}{Aug 2024}
    \item Prototyped diffusion-based and VAE-based generative models to synthesize training data for low-resource conditions; developed a real-time multimodal detection system over combined audio and sensor streams.
\end{cvrole}

\begin{cvrole}{Visiting Researcher}{University of Visual and Performing Arts}{Colombo 02, Sri Lanka}{Jan 2024}{Mar 2024}
    \item Ran human-centred evaluations of interactive machine learning systems, pairing quantitative metrics with structured operator feedback to produce usability and adoption recommendations.
\end{cvrole}

\begin{cvrole}{Technical Consultant}{Forestpin (Pvt) Ltd}{Colombo 05, Sri Lanka}{Aug 2020}{Feb 2021}
    \item Built machine learning and statistical modelling pipelines for financial forensics, compliance monitoring, and risk detection over large enterprise datasets.
    \item Developed SQL and Python backend scoring services that flagged anomalous transactions for analyst review.
    % , with human confirmation required before any action was taken on a flagged record.
    % \item Worked directly with compliance and risk stakeholders to translate detection requirements into technical specifications and acceptance criteria.
\end{cvrole}

\begin{cvrole}{Research Engineer}{MAS Holdings (Pvt) Ltd}{Colombo, Sri Lanka}{Jan 2015}{Jul 2020}
    \item Owned the end-to-end development and deployment of machine learning and computer vision systems for quality assurance in industrial manufacturing settings; improved operational efficiency by up to 300\%.
    \item Built a document digitization pipeline combining OCR and layout segmentation to convert a large physical archive into a searchable corpus; delivered a context-aware search system integrated into the company ERP, still in use today.
    \item Developed real-time C++ and Python inference services on embedded hardware, and built ingestion and ETL pipelines for high-frequency operational IoT data across distributed sites.
    \item Deployed human-in-the-loop inspection models that balanced throughput against operational risk; mentored engineers and established code review and CI/CD practice to scale adoption across teams.
\end{cvrole}

% =====================================================
\section*{Benchmark Datasets \footnotesize {\normalfont{(published on Zenodo; 7,000+ recordings from 70 contributors)}}}
\begin{itemize}[leftmargin=1.2em, itemsep=-0.25em]
    % \item \textbf{DoMP} (4,000 symbolic sequences, \href{https://doi.org/10.5281/zenodo.10818617}{DOI 10818617}), \textbf{DoPP} (2,000 raw audio, \href{https://doi.org/10.5281/zenodo.14497998}{DOI 14497998}), \textbf{DoDP} and \textbf{DoDP2} (2,000 each, \href{https://doi.org/10.5281/zenodo.14497974}{DOI 14497974}, \href{https://doi.org/10.5281/zenodo.18395007}{DOI 18395007}). Designed with annotation protocols and held-out evaluation splits for cross-group comparability.

    \item \textbf{DoMP}, 4,000 recordings, symbolic monophonic patterns with expressive variations (\textit{DOI: \href{https://doi.org/10.5281/zenodo.10818617}{\textit{10818617}}})
    \item \textbf{DoPP}, 2,000 recordings, polyphonic audio patterns with expressive variations (\textit{DOI: \href{https://doi.org/10.5281/zenodo.14497998}{\textit{14497998}}})
    \item \textbf{DoDP}, 2,000 recordings, percussive audio, rhythmic pattern detection (\textit{DOI: \href{https://doi.org/10.5281/zenodo.14497974}{\textit{14497974}}})
    \item \textbf{DoDP2}, 2,000 recordings, percussive audio, extended ontology and multi-performer coverage (\textit{DOI: \href{https://doi.org/10.5281/zenodo.18395007}{\textit{18395007}}})
\end{itemize}

% =====================================================
\section*{Selected Publications \footnotesize {\normalfont{(full list at: \href{https://nishal.xyz/publications}{\textit{nishal.xyz/publications}})}}}
\begin{itemize}[leftmargin=1.2em, itemsep=-0.25em]
    \item Silva, N. and Turchet, L. \textit{Real-Time Audio Pattern Detection for Smart Musical Instruments}. Journal of the Audio Engineering Society, Mar. 2026. 
    \item Silva, N., Wanderley, M. and Turchet, L. \textit{Melody and Motion: Integrating Guitar Gesture Detection with Musical Patterns for Extended Control..}. IEEE International Symposium on the Internet of Sounds, Oct. 2025.
    \item Silva, N., Boem, A. and Turchet, L. \textit{Interactive IoMusT-Based Concerts: Real-Time Pattern Recognition and Audience Experience}. IEEE International Conference on Immersive and 3D Audio, Sep. 2025.
\end{itemize}

% =====================================================
\section*{Selected Projects \footnotesize {\normalfont{(full portfolio at: \href{https://nishal.xyz/research}{\textit{nishal.xyz/research}})}}}
\begin{itemize}[leftmargin=1.2em, itemsep=-0.25em]
    \item \textbf{Distributed Multi-Agent Coordination Network:} Event-driven architecture where on-device inference models routed decisions to networked peripherals over custom protocols. Architecture supported one-to-one, one-to-many, and many-to-many mappings under hard latency budgets. Published at IEEE I3DA 2025.
    \item \textbf{Archive Digitization and Context-Aware Search:} OCR and layout segmentation pipeline converting a physical archive of technical documents into a searchable corpus, with a retrieval interface embedded in an enterprise ERP.
    \item \textbf{Synthetic Data Generation for Low-Resource Training:} Conducted extensive research on rule-based, diffusion-based and VAE-based generation schemes to address enhanced augmentation in low training data availability.
\end{itemize}

% =====================================================
\section*{Service, Recognition \& Highlights}
\begin{itemize}[leftmargin=1.2em, itemsep=-0.25em]
    \item \textbf{Peer Review:} IEEE Access, Journal of the National Science Foundation (Sri Lanka). \textbf{Program Committee:} IEEE Internet of Sounds Symposium, IEEE Immersive and 3D Audio Conference.
    \item \textbf{Finalist, MIDI Innovation Awards} (Sep 2023), real-time symbolic pattern detection for smart instruments. 
    \item \textbf{GMOA Recognition (Sri Lanka)} for a computer vision system monitoring COVID-19 patient vitals across hospital wards, reducing healthcare worker exposure.
\end{itemize}

\end{document}